\section{INTRODUCTION}
\label{sec:introduction}

\subsection{Background and Motivation}

Communication and sensing are often needed within the same system. A mobile
terminal needs a data link and knowledge of its position, while a network
operator needs to carry traffic and monitor the infrastructure that supports
it. Integrated sensing and communication (ISAC) brings these functions together
through shared signals, hardware, or infrastructure
\cite{Zhang2026ISACEvolution}. This interest extends to future mobile networks,
with ISAC included among the usage scenarios in the International Mobile
Telecommunications framework for 2030 and beyond
\cite{ITUR_M2160_2023}.

Optical integrated sensing and communication (O-ISAC) applies this idea to
optical links and photonic systems. A fiber can carry communication and sensing
signals along the same route, while a visible light source can support data
transmission and positioning \cite{OISAC_SCR00013,OISAC_SCR00196}.
Figure~\ref{fig:oisac_platform_overview} illustrates representative ways of
combining the two functions. The role of optics differs across these systems.
In fiber and optical wireless links, light propagates through the sensing
medium or reaches the receiver or target. In photonics-enabled wireless
systems, optical components can generate the signal while propagation to the
target takes place in the millimeter-wave or terahertz band
\cite{OISAC_SCR00036,OISAC_SCR00083}.

\begin{figure*}[!t]
\centering
\includegraphics[width=\textwidth]{figures/fig_introduction_oisac_platforms.png}
\caption{Representative O-ISAC configurations in fiber, free-space optical,
visible-light, and photonics-enabled terahertz systems. Each panel shows a
shared resource and the communication and sensing functions it supports.
The design choices below the panels can affect both functions, with their
roles depending on the implementation. These examples illustrate physical
configurations rather than a performance ranking or an exhaustive taxonomy.}
\label{fig:oisac_platform_overview}
\end{figure*}

Sharing a resource can also link the performance of the two functions.
In one fiber experiment, raising the sensing-probe power while holding
communication launch power fixed improved the sensing signal-to-noise ratio
but increased the communication bit error rate \cite{OISAC_SCR00057}.
Joint operation can also allow sensing to assist communication. In a
visible-light prototype, position estimates guided beamforming toward the
receiver and improved the communication data rate at the tested locations
\cite{OISAC_SCR00196}. These examples show why the shared resource and the way
it is used matter. Joint design must account for both the benefit sought and
its effect on the other function.

The physical setting determines how these relationships can be used.
A vibration measurement along a fiber and a position estimate in a room
answer different sensing questions. Optical wireless links also depend on
propagation and alignment, while fiber systems have their own transmission
and sensing constraints \cite{OISAC_SCR00916,OISAC_SCR01085,Mohsan2026OISACReview}.
A researcher selecting an O-ISAC approach therefore needs more than a list
of achievable rates or sensing errors. The useful questions are which design
choice produced the result, how both functions responded, and whether the
conditions match the intended application. This connects the reported
performance to decisions about how to build and evaluate a system.

\subsection{Related Surveys}
\label{sec:related_surveys}

Existing surveys explain O-ISAC from several complementary perspectives.
Architecture-oriented works discuss how communication and sensing can be
combined, while studies of analytical limits and estimation explain the
performance available under particular models
\cite{Wen2024OISACArchitectures,Su2025OISACTheoryPractice,
Khorasgani2026GeneralOpticalISACSurvey}. More focused reviews examine
visible-light integration, converged fiber networks, and photonic waveform
design \cite{Liang2024LiSACReview,Liu2026OpticalFiberISACReview,
Lu2026ConvergedFiberISACReview,Lyu2026PhotonicTHzWaveforms}.
Table~\ref{tab:prior_surveys} compares representative works along six dimensions,
from physical scope to research directions. Each entry identifies how the
work treats that dimension.

\begin{table*}[!t]
\caption{Comparison of representative related reviews and tutorial-style
syntheses with the present survey}
\label{tab:prior_surveys}
\centering
\begingroup
\footnotesize
\setlength{\tabcolsep}{3pt}
\renewcommand{\arraystretch}{1.16}
\begin{tabularx}{\textwidth}{@{}%
>{\raggedright\arraybackslash}p{0.105\textwidth}
>{\raggedright\arraybackslash}X
>{\raggedright\arraybackslash}X
>{\raggedright\arraybackslash}X
>{\raggedright\arraybackslash}X
>{\raggedright\arraybackslash}X
>{\raggedright\arraybackslash}X@{}}
\toprule
\textbf{Work} &
\textbf{Optical scope} &
\textbf{Shared resources} &
\textbf{Numerical results} &
\textbf{Joint design} &
\textbf{Conditions and validation} &
\textbf{Research directions} \\
\midrule
Wen et al.\ \cite{Wen2024OISACArchitectures}\newline 2024 &
Free-space optical links &
Optical front end and waveform &
Bit-error and ranging-error simulations &
Bias, power, and subcarrier allocation &
Waveform, link, and channel settings &
Hardware models and RF--optical coordination \\
\addlinespace[3pt]
Liang et al.\ \cite{Liang2024LiSACReview}\newline 2024 &
Visible light and hybrid links &
Lighting, positioning signals, and receivers &
Communication and sensing performance tables &
Lighting constraints and positioning-assisted links &
Prototype indicators and system-specific tests &
Blockage, uplinks, and hybrid integration \\
\addlinespace[3pt]
Lyu et al.\ \cite{Lyu2026PhotonicTHzWaveforms}\newline 2026 &
Photonic millimeter-wave and terahertz systems &
Integrated waveforms and photonic front ends &
Demonstration table and experimental comparisons &
Weighted objectives and waveform parameters &
Carrier, path, power, and receiver settings &
Stable sources, beam steering, and security \\
\addlinespace[3pt]
Zhang et al.\ \cite{Zhang2026ISACEvolution}\newline 2026 &
Optical and photonic systems within wider ISAC &
Cooperative and shared-signal architectures &
Optical performance examples &
RF--optical resource tradeoffs &
Alignment, timing, and testbeds &
Adaptive task allocation and joint timing \\
\addlinespace[3pt]
Mohsan et al.\ \cite{Mohsan2026OISACReview}\newline 2026 &
Fiber, optical wireless, and photonic wireless &
Shared fiber and joint waveforms &
Data-rate and sensing-result tables &
Waveform design and resource allocation &
Experimental setups and test spaces &
Predistortion, calibration, and 3D positioning \\
\midrule
\textbf{This survey} &
Fiber, optical wireless, photonic wireless, and hybrids &
Physical paths and sharing mechanisms &
Metric definitions and source-specific interpretation &
Mechanisms linking both functional outcomes &
Joint test conditions and reproducibility &
Benchmarks and paired disturbance tests \\
\bottomrule
\end{tabularx}
\par\vspace{4pt}
\noindent
\parbox{\textwidth}{\raggedright\footnotesize
The entries describe each work's analytical approach within its stated scope.
For the broader ISAC survey by Zhang et al., the entries concern its optical
and RF--optical treatment. Numerical coverage includes worked examples,
reported demonstrations, or interpretation of reported metrics, as specified
in each cell.\par}
\endgroup
\end{table*}

Recent surveys also connect several parts of the field. Zhang et al.
\cite{Zhang2026ISACEvolution} examine the evolution of ISAC across spectrum,
network architecture, and sensing modalities, together with security and
standardization. Mohsan et al. \cite{Mohsan2026OISACReview} review optical
platforms and experimental demonstrations alongside waveform design,
applications, and open challenges. Numerical demonstrations, joint design,
and research directions are therefore already important parts of the survey
literature. Our contribution is to connect these elements within a systematic
review organized around shared resources. We examine how a design choice
affects both functions, which measurement and operating conditions determine
its interpretation, and what joint tests would support its use in another
system.

\subsection{Scope and Contributions}

This survey examines 206 unique O-ISAC studies first made publicly available
from 1 January 2020 through 22 June 2026. It covers guided optical systems,
optical wireless links, and photonics-enabled wireless systems, including
hybrid configurations. Across these settings, we examine how communication
and sensing are connected through a physical resource, a signal, or a joint
design objective. The review follows three questions. What is shared between
the two functions? How does the shared design affect their performance?
What further evaluation is needed to use the resulting insight in another
system?

The survey makes three contributions.
\begin{itemize}
\item We distinguish optical platforms by their physical signal paths and
identify how shared hardware, signals, and resource allocation connect
communication and sensing. This gives the reader a basis for selecting
integration mechanisms that fit the intended physical setting.

\item We synthesize the mechanisms through which bandwidth, power, and
waveform or receiver design affect communication rate and reliability
alongside sensing accuracy and resolution. We interpret reported outcomes
using their metric definitions and operating conditions, so that a design
benefit can be read alongside its effects on the other function.

\item We connect validation evidence to concrete evaluation priorities,
including joint tests under traffic, changing illumination, atmospheric
variation, and optical drift. These priorities identify the need for paired
performance measurements, separate-function baselines, and reproducible
configurations when testing a design beyond its original setting.
\end{itemize}

Section~\ref{sec:foundations_comparison} introduces the physical and performance
concepts used in the review, and Section~\ref{sec:review_methods} describes the
review methods. Sections~\ref{sec:optical_modality_families}
through~\ref{sec:technologies_applications_6g} examine platforms, joint design,
validation, and applications. Section~\ref{sec:discussion_roadmap} develops the
research priorities, and Section~\ref{sec:conclusion} concludes the article.

\FloatBarrier
